\begin{textsample}{POS Dim 2 – ap – Score 23.00 – ap/ap\_em\_42.txt}  \label{ex:f2_pos_146}
A \textbf{Portaria} 723/2012 criou o Cadastro Nacional de Aprendizagem Profissional - CNAP, no qual devem se inscrever todas as \textbf{entidades} qualificadas em formação técnico-profissional metódica.
Para as \textbf{entidades} sem fins lucrativos dedicadas à educação \textbf{profissional}, a inserção no CNAP dependerá de avaliação de competência, com vistas à verificação de sua aptidão para ministrar programas que permitam a inclusão de aprendizes no mercado de trabalho.
No Novo Ensino Médio a formação \textbf{técnica} e \textbf{profissional} passa a fazer parte do Ensino Médio regular.
Isso quer dizer que mesmo estudantes que não escolherem estudar em uma escola \textbf{técnica} no início da etapa podem escolher compor parte ou toda a sua \textbf{carga} \textbf{horária} \textbf{destinada} aos \textbf{itinerários} com \textbf{cursos} \textbf{técnicos} ou FIC, a partir da disponibilidade de \textbf{oferta} em seu território.
A flexibilidade dos \textbf{perfis} \textbf{profissionais}, a permeabilidade de funções entre ocupações diferentes, o desaparecimento de ocupações e o surgimento de outras em grande velocidade são desafios para a organização do \textbf{currículo} e para a criação de propostas educativas atraentes, dinâmicas e aderentes à realidade.
Esse cenário complexo e instável traz um duplo desafio : desenvolver \textbf{itinerários} \textbf{formativos} condizentes com os \textbf{itinerários} \textbf{profissionais} \textbf{vigentes} e elaborar \textbf{currículos} com coerência interna e convergentes entre si.
Dessa forma, justifica -se o esforço dedicado para a estruturação de \textbf{itinerários} \textbf{formativos} que \textbf{atendam} às atuais necessidades do mundo do trabalho e às expectativas dos alunos.
O Modelo Pedagógico da SEED, nesse sentido, ao apresentar os princípios e concepções educacionais, ao explicitar a lógica curricular dos \textbf{cursos} de Educação Profissional e ao orientar a prática educativa desenvolvida na Instituição, assume a importância da organização do portfólio de seus \textbf{cursos} presenciais e a distância em \textbf{itinerários} \textbf{formativos} \textbf{atualizados}.
A incorporação cada vez mais intensa das tecnologias na rotina e nos processos organizacionais tem trazido mudanças significativas na dinâmica do trabalho.
Conectividade, inovação e automação são algumas das transformações que definem os novos contornos da economia global, afetando diretamente a organização social e \textbf{técnica} do trabalho.
Esse cenário requer da educação \textbf{profissional} um modelo de formação amplo, capaz de articular conhecimentos, habilidades, atitudes e valores para o desenvolvimento de competências que possam ser mobilizadas e aprimoradas ao longo de toda a vida.
A adoção de metodologias ativas e a ênfase no protagonismo do aluno no decorrer do processo de ensino e aprendizagem contribuem para a formação de pessoas capazes de lidar com processos complexos e preparadas para \textbf{atender} às múltiplas exigências do segmento de atuação.
Para além do domínio técnico-científico próprio de cada ocupação, é preciso que a educação \textbf{profissional} promova o \textbf{empreendedorismo}, a sustentabilidade e a laboralidade, numa perspectiva crítica e comprometida com transformação da realidade.
É nessa conjuntura que se efetiva a aprendizagem por competência, a possibilidade de aproveitamento de estudos e da experiência \textbf{profissional} do \textbf{trabalhador} e, principalmente, a \textbf{oferta} de \textbf{cursos} de educação \textbf{profissional} elaborados com base em \textbf{itinerários} \textbf{formativos}, de forma a \textbf{alinhar} as expectativas dos alunos, as demandas das empresas e a \textbf{oferta} de educação \textbf{profissional}.
Considerando esses fatores, a organização de \textbf{itinerários} deve se pautar numa perspectiva de educação contínua, para a vida e por toda a vida, já que tem como \textbf{prerrogativas} aumentar as chances de os alunos assumirem o próprio percurso \textbf{formativo} e consolidar a ideia de que a educação \textbf{profissional} favorece a inserção das pessoas no mundo do trabalho.
A concepção de \textbf{itinerários} \textbf{formativos} utiliza como base a \textbf{legislação} educacional \textbf{vigente}.
Apesar de o tema já ter sido tratado nos pareceres CNE/CEB nº 16/1999 e CNE/CEB nº 11/2012, é com a publicação do Decreto Federal nº 5.154, de 23 de julho de 2004 que se \textbf{normatiza} a organização da educação \textbf{profissional} em \textbf{itinerários} \textbf{formativos}, ainda que essa expressão não seja utilizada na redação do decreto.

% matched lemmas: alinhar, atender, atualizar, carga, currículo, curso, destinar, empreendedorismo, entidade, formativo, horário, itinerário, legislação, normatizar, oferta, perfil, portaria, prerrogativa, profissional, trabalhador, técnico, vigente
\end{textsample}
